% ============================================================
%  FTN-LoRA: prompt-perplexity-routed LoRA atom bank
%  Self-contained Overleaf main. Compile with pdfLaTeX.
%
%  Transcribed from the implementation:
%    model.py / lora_bank.py / perplexity_configurer.py
%    / spatial.py / train.py
% ============================================================
\documentclass[11pt]{article}

\usepackage[utf8]{inputenc}
\usepackage[T1]{fontenc}
\usepackage{lmodern}
\usepackage[margin=1in]{geometry}
\usepackage{amsmath,amssymb,amsthm}
\usepackage{mathtools}
\usepackage{bm}
\usepackage{microtype}
\usepackage{hyperref}
\hypersetup{colorlinks=true,linkcolor=blue,citecolor=blue,urlcolor=blue}

% ---- Math operators / shorthands ---------------------------
\DeclareMathOperator*{\argmin}{arg\,min}
\DeclareMathOperator{\Topk}{Top\text{-}k}
\newcommand{\kwta}{\mathrm{kWTA}_k}
\newcommand{\sg}{\mathrm{sg}}        % stop-gradient
\newcommand{\drop}{\mathrm{drop}}
\newcommand{\vecop}{\mathrm{vec}}
\newcommand{\grid}{\mathrm{grid}}
\newcommand{\CE}{\mathrm{CE}}
\newcommand{\RR}{\mathbb{R}}
\newcommand{\Norm}[1]{\mathcal{N}\!\left(#1\right)}

\title{FTN-LoRA: A Prompt-Perplexity-Routed LoRA Atom Bank}
\author{}
\date{}

\begin{document}
\maketitle

\begin{abstract}
We attach a bank of $H$ low-rank (LoRA) adapter \emph{atoms} to every target
linear of a frozen language-model backbone. A single shared binary mask selects
$k$ active atoms; at inference the mask is chosen \emph{self-supervisedly} by
minimising the next-token perplexity of the prompt over the mask, with all
weights frozen. Tasks are learned block-sequentially: only the atoms selected
for a task are updated, and a mask-overlap-filtered replay protects earlier
blocks. The equations below transcribe the reference implementation.
\end{abstract}

% ---- Notation ----------------------------------------------
\section{Notation}
The frozen backbone is $\phi$ (Qwen2.5-7B-Instruct) with parameters $\theta$ held
fixed. A bank is attached to every target linear
$\ell \in \{\texttt{q},\texttt{k},\texttt{v},\texttt{o},\texttt{up},\texttt{gate},\texttt{down}\}\texttt{\_proj}$.
There are $H = g^2$ atoms laid out on a $g \times g$ grid (default $g = 48$, so
$H = 2304$); $k$ atoms are active per forward pass (default $k = 128$); each atom
has rank $r$ (default $r = 1$) and the bank delta is scaled by $\alpha$ (default
$\alpha = 16$).

% ---- 1. Banked-LoRA layer ----------------------------------
\section{Banked-LoRA layer}
\label{sec:layer}
For a target linear $\ell$ with frozen weight $W_\ell \in \RR^{\text{out}\times\text{in}}$,
the bank stores $A_\ell \in \RR^{H\times\text{in}\times r}$ and
$B_\ell \in \RR^{H\times r\times\text{out}}$. With a shared binary mask
$m\in\{0,1\}^H$ (where $\|m\|_1 = k$),
\begin{equation}
  y_\ell \;=\; W_\ell\, x \;+\; \frac{\alpha}{k}
  \sum_{j=1}^{H} m_j \,\bigl(\,\drop(x)^{\!\top} A_{\ell,j}\,\bigr)\, B_{\ell,j}.
\end{equation}
Equivalently, with the live active set $S(m) = \{\, j : m_j = 1 \,\}$,
\begin{equation}
  y_\ell \;=\; W_\ell\, x \;+\; \frac{\alpha}{k}
  \sum_{j \in S(m)} \bigl(\drop(x)^{\!\top} A_{\ell,j}\bigr) B_{\ell,j}.
\end{equation}
A single mask $m$ is shared across all layers $\ell$, so one selection drives the
whole model. The scale $\alpha/k$ tracks the \emph{live} active-atom count, so $k$
may be changed at inference without inflating or fading the delta.

% ---- 2. kWTA ----------------------------------------------
\section{$k$-winners-take-all}
\begin{equation}
  \bigl[\kwta(z)\bigr]_j \;=\;
  \begin{cases}
    1, & j \in \Topk(z), \\
    0, & \text{otherwise.}
  \end{cases}
\end{equation}

% ---- 3. Prompt-perplexity reconfiguration ------------------
\section{Prompt-perplexity reconfiguration}
\label{sec:reconfig}
This is the core routing mechanism. Given prompt token-ids
$u = (u_1,\dots,u_T)$, we optimise mask logits $z \in \RR^H$ to minimise the
self-supervised next-token NLL, with $\theta$ and $\{A,B\}$ frozen:
\begin{equation}
  \mathcal{L}_{\mathrm{ppl}}(m) \;=\;
  -\frac{1}{T}\sum_{t=1}^{T} \log p_{\theta,\,m}\!\bigl(u_t \mid u_{<t}\bigr).
\end{equation}
The AnnealSTE inner loop runs for steps $s = 0,\dots,N-1$; only $z$ receives
gradient:
\begin{align}
  \tau_s &= \operatorname{std}(z)\cdot \exp\!\Bigl(-\,s \big/ \tfrac{N}{4}\Bigr),
            &\varepsilon_s &\sim \Norm{0,\tau_s^2 I}, \\
  m_s    &= \kwta(z + \varepsilon_s) \;+\; \sigma(z) - \sg\!\bigl(\sigma(z)\bigr),
            &&\text{(hard forward, straight-through grad)}\\
  z      &\leftarrow z - \eta\,\mathrm{Adam}\!\bigl(\nabla_z \mathcal{L}_{\mathrm{ppl}}(m_s)\bigr).
\end{align}
The best hard mask is retained by NLL:
\begin{equation}
  m^\star = \kwta(z_{s^\star}), \qquad
  s^\star = \argmin_{s} \mathcal{L}_{\mathrm{ppl}}\!\bigl(\kwta(z_s)\bigr).
\end{equation}
An optional warm start (e.g.\ mid-rollout reconfiguration) initialises
$z_0 = a\, m_{\text{init}} + b\,\xi$ with $\xi \sim \Norm{0,I}$.

% ---- 4. Lateral smoothing ----------------------------------
\section{Lateral smoothing}
The atoms are reshaped to the $g\times g$ grid; $C$ is a fixed $3\times3$ box
convolution (circular padding iff the grid is a torus). We iterate $n_{\mathrm{lat}}$
Polyak-averaged smoothing passes and re-apply $\kwta$ to pull active atoms into
contiguous blobs:
\begin{align}
  m &\leftarrow \rho\, m + (1-\rho)\,\vecop\!\bigl(C * \grid(m)\bigr),
     \qquad t = 1,\dots,n_{\mathrm{lat}}, \\
  m_{\mathrm{final}} &= \kwta(m).
\end{align}

% ---- 5. Per-task reference mask ----------------------------
\section{Per-task reference mask}
The reference mask for task $c$ accumulates the masks over its support
reconfigurations, then takes $\kwta$; mask overlap is measured by
\begin{equation}
  m^{\mathrm{ref}}_c \;=\; \kwta\!\Bigl( \textstyle\sum_{e \in \mathrm{supp}(c)} m_e \Bigr),
  \qquad
  o(m, m') = \frac{\langle m, m'\rangle}{k} \in [0,1].
\end{equation}

% ---- 6. Continual training objective -----------------------
\section{Continual training objective}
Training is block-sequential. For task $c$ we freeze $\theta$; for each batch
$B_c$ we first reconfigure the mask $m$ on one random member's prompt
(Section~\ref{sec:reconfig}), then update \emph{only} the active atoms on the
causal-LM loss. Because $m_j = 0$ zeroes atom $j$'s contribution in
Section~\ref{sec:layer}, its gradient is $0$ and untouched atoms are protected:
\begin{equation}
  \min_{A,B}\;
  \underbrace{\frac{1}{|B_c|}\sum_{(u,y)\in B_c} \CE\bigl(y, p_{\theta,m}(\cdot\mid u)\bigr)}_{\text{current task, under fixed }m}
  \;+\;
  \mathbb{1}[r<p_{\mathrm{rep}}]\cdot
  \underbrace{\CE\bigl(y', p_{\theta,m'}(\cdot\mid u')\bigr)}_{\text{replayed prior example, under its stored }m'} .
\end{equation}
Replay candidates are filtered to mid-band mask overlap with the current
selection --- the retention mechanism:
\begin{equation}
  (u', y', m')\ \text{drawn from prior buffers with}\quad
  0.1 \;<\; o(m', m) \;\le\; 0.8 .
\end{equation}

\end{document}
